\section{热力学第一定律}
\label{sec:热力学第一定律}
\subsection{热力学过程和平衡过程}
\label{subsec:热力学过程和平衡过程}
热力学系统的状态随时间发生变化的过程称为\uwave{热力学过程}，简称为过程，在热力学过程的进行的任一时刻，系统的状态一般不是平衡态。例如当我们推动活塞压缩气缸内气体时，此时气缸内的温度和压强的分布未必是均匀的，因此我们就不能简单的用一组状态量$(P,V,T)$来描述气缸内理想气体的性质，简单的说，\empx{热力学过程持续改变理想气体的稳态，而达到稳态需要弛豫时间，这就导致过程中任意一个时刻都不是稳态}。这就为研究热力学过程带来了麻烦。

但是，如果实际的物理过程的进行速度非常缓慢，以至于系统的弛豫时间远小于系统状态发生每一个微小变化的时间，那么此时系统就总是能及时的响应外界条件的改变并立即达到新的稳态。不妨仍然以推动活塞压缩气缸内气体为例，由于活塞的推动使得气体各处密度不同，则会中作用在气体中以疏密波的形式传播，其传播速度为气体声速，约在$10^2\si{m\cdot s^{-1}}$的量级，因此对于线度约为$1\si{m}$的容器而言，压强恢复均匀一致所需的时间，或者说，压强的弛豫时间就大约为$10^{-2}\si{s}$，这样只要控制推动活塞的速度，使得$10^{-2}\si{s}$内活塞没有显著的位置改变，那么这就足矣使得容器内的气体及时的响应体积的压缩，跟随变化实时达到平衡态，这样我们就达成了一种平衡态的连续变化，这种想法抽象出来，就形成了\uwave{准静态过程}的想法，即\empx{热力学过程进行的足够缓慢，以至于过程连续经历的每一个中间状态都可以视为处于平衡态}。在接下来的内容中，我们总是会考虑平衡过程，这是一种非常有价值的简化，这样就可以用\xref{chap:气体动理论}中研究平衡态的方法研究平衡过程中的每一状态，从而回避变化过程中繁琐且不必要的微分方程。

\subsection{热力学系统中的能量变化}
\label{subsec:热力学系统中的能量变化}
\subsubsection{准静态过程中的功}
力学研究表明，外力对系统做功将改变系统的状态，并且伴有以功为形式的能量转移。这一点在热力学系统中也是成立的，当热力学系统与外界存在能量交换时，做功也是系统能量转移的一种方式。在这里我们只考虑准静态过程中理想气体的体积发生变化时，理想气体的压力对外界所做的机械功，也称为\uwave{体积功}。如\xref{fig:准静态过程的功}所示，设想气缸内的气体经历了无摩擦的准静态膨胀过程，推动气缸的活塞外移，对外做功。记气体压强为$P$，记录气缸活塞的面积为$S$。

\begin{OldFigure}[准静态过程的功]
    \inputfigure[scale=0.9]{Chapter12A_Figure01}
\end{OldFigure}

现考虑当活塞外移$\dd{l}$时气体对外所做的功$\dd{W'}$，由于气体在活塞上的作用力为$F=PS$
\begin{Equation}
    \dd{W'}=PS\dd{l}
\end{Equation}
这里我们注意到气体的体积增量为$\dd{V}=S\dd{l}$，故上式可以写作
\begin{Equation}
    \dd{W'}=P\dd{V}
\end{Equation}
这里计算的是气体对外做功，但是如果从系统的视角出发，我们将采取不同的正负记法
\begin{itemize}
    \item 记外界对系统做功的$W$为正，因为这使得系统的能量增加，该过程中气体压缩。
    \item 记系统对外界做功的$W$为负，因为这使得系统的能量减少，该过程中气体膨胀。
\end{itemize}
由于$W$的正值被定为外界对系统的做功，因此$W$的公式还要在上述基础上取负。\footnote{这里采取的功和热量的正值记法都是以使得系统能量增加为标准的，与包括\cite{b2a}和\cite{b3a}在内的主流大学物理教材和\cite{b4}的大学化学教材是一致的，然而由上海大学编写的教材\cite{b1}中却记功的正值方向是“系统对外界做功”而非自然的“外界对系统做功”，如果需要顺从它的这种有些怪异的记法，因在所有出现$W$的地方用$-W$代替，并且在功$W$的计算公式中就不再需要出现负号了。}
\begin{OldBoxFormula}[理想气体的体积功]*
    理想气体的体积功可以表示为，其中$W$表示外界对气体做功
    \begin{Equation}&[]
        W=-\Int[V_1][V_2]P(V)\dd{V}
    \end{Equation}
    相应的微元形式是\footnote[2]{这里我们注意到微元形式中的功使用$\fdd{W}$而不是通常的$\dd{W}$表示，稍后我们会解释这样写的意义。}
    \begin{Equation}&[微元]
        \fdd{W}=-P\dd{V}
    \end{Equation}
\end{OldBoxFormula}

我们知道，理想气体的准静态过程可以用P--V图上的曲线表示，而\fancyref{fml:理想气体的体积功}告诉我们，在P--V图上的曲线下面积表示体积功的大小，如\xref{fig:理想气体的体积功的图示}所示
\begin{OldFigure}[理想气体的体积功的图示]
    \begin{OldFigureSub}[压缩过程]
        \inputfigure[scale=0.9]{Chapter12A_Figure03}
    \end{OldFigureSub}
    \hspace{1cm}
    \begin{OldFigureSub}[膨胀过程]
        \inputfigure[scale=0.9]{Chapter12A_Figure02}
    \end{OldFigureSub}
\end{OldFigure}

特别需要注意的是正负的问题，\empx{压缩过程中功为正，膨胀过程中功为负}。这很容易记忆，因为通常我们认为自左向右的积分将给出正的面积，而这里的功是积分的负值，即曲线下面积的负值，所以说，自右向左的压缩P--V曲线将给出正功，自左向右的膨胀P--V曲线将给出负功。

\subsubsection{准静态过程中的热}
外界对系统做功使得系统的状态发生变化，这是热力学系统能量交换的一种方式，然而，这并不是热力学系统能量交换的唯一方式，试想，两个温度不同的物体具有不同的内能，而将两者接触后，两者的温度将达到平衡，即内能相同，因而这个过程中必然发生了能量的转移，但很明显没有任何的做功，因此做功之外的能量转移方式必然是存在的。我们将因系统间的温度差而在系统之间传递热运动能量的方式，称为\uwave{热传导}，这种传递的热运动能量称为\uwave{热量}，通常我们用符号$Q$表示，由于热量的本质仍然是能量转移的多少，因此热量的单位同样是焦耳。

热量在焦耳测定热功当量之前的时代中，还使用过称为卡路里的单位
\begin{Equation}*
    1\si{cal}=4.181\si{J}
\end{Equation}
这里卡路里\si{cal}的物理意义是将$1\si{g}$水在$1\si{atm}$的气压下提升$1\si{\degreeCelsius}$所需吸收的热量，目前在物理学领域中卡路里已经很少使用，但是在营养标签中还常用千卡$\si{kcal}$作为食物热量的标识。

热量和功在本质上是完全两种不同的概念，做功是通过系统的宏观位移完成的，\empx{是宏观的机械运动与微观的无规则热运动之间的能量转化}，例如压缩气缸内的气体就是将推动活塞的机械运动转化为了气体分子的热运动。热传递则是通过微观分子间的相互作用完成的，两个温度不同的系统就是两部分平均速率不相同的分子，两者的温度达到平衡的过程，实质上就是两者的分子通过碰撞实现速率一致的过程，因此热传导\empx{是微观的无规则热运动内部的能量转移}。

热量和功都是对系统能量变化的量度，两者的共同特征是，两者均是与变化过程有关的量。

热量的正负记法遵循与做功类似的规则
\begin{itemize}
    \item 记外界对系统传热的$Q$为正，因为这使得系统的能量增加，该过程中气体吸热。
    \item 记系统对外界传热的$Q$为负，因为这使得系统的能量减少，该过程中气体放热。
\end{itemize}

\subsection{热力学第一定律}
\label{subsec:热力学第一定律}
通过\xref{subsec:热力学系统中的能量变化}，我们知道热力学系统可以通过力学作用“做功”和热学作用“传热”这两种形式与外界发生能量交换，而另外一方面，我们知道热力学系统自身的能量是内能，由此依据能量守恒定律，就可以建立“~\textbf{功$W$}、\textbf{热$Q$}、\textbf{内能$E$}~”间的关系，即著名的热力学第一定律。
\begin{OldBoxLaw}[热力学第一定律]
    热力学系统内能的增量，等同于外界对系统所做功与外界传递给系统的热量之和
    \begin{Equation}&[]
        \delt{E}=Q+W
    \end{Equation}
    而对于其中的微元过程，上式可以改写为
    \begin{Equation}&[微元]
        \dd{E}=\fdd{Q}+\fdd{W}
    \end{Equation}
    该结论称为\uwave{热力学第一定律}。
\end{OldBoxLaw}

关于热力学第一定律中各个符号正负的意义，我们总结在\xref{tab:热力学第一定律中的物理量}
\begin{OldTable}[热力学第一定律中的物理量]{llc<{$\ >0$}llc<{$\ <0$}ll}
<\mc{1}{物理量}&类型&\mc{3}{正值含义}&\mc{3}{负值含义}\\>
功&过程量&$W$&气体压缩&外界对系统做功&$W$&气体膨胀&系统对外界做功\\
热&过程量&$Q$&气体吸热&外界对系统传热&$Q$&气体放热&系统对外界传热\\
内能&状态量&$E$&\mc{2}{气体从外界获取能量}&$E$&\mc{2}{气体向外界释放能量}\\
\end{OldTable}
关于热力学第一定律中的\xrefpeq{微元}，注意到使用了$\fdd{Q},\fdd{W}$而不是更常见的$\dd{Q},\dd{W}$，这其实是因为内能$E$是关于$(P,V,T)$的状态量，而$Q$和$W$是过程量，无法定义通常意义上的微分。

接下来我们从数学上来说明这个问题，根据\fancyref{fml:理想气体的内能}
\begin{Equation}*
    E=\frac{i}{2}\nu RT
\end{Equation}
这样内能$E$的全微分$\dd{E}$就可以表示为
\begin{Equation}[内能的微分]
    \dd{E}=\pdv{E}{P}\dd{P}+\pdv{E}{V}\dd{V}+\pdv{E}{T}\dd{T}=\frac{i}{2}\nu R\dd{T}
\end{Equation}
根据\fancyref{fml:理想气体的体积功}
\begin{Equation}[功的微分]
    \fdd{W}=-P\dd{V}
\end{Equation}
根据\fancyref{law:热力学第一定律}和上两式
\begin{Equation}[热的微分]
    \fdd{Q}=\dd{E}-\fdd{W}=P\dd{V}+\frac{i}{2}\nu R\dd{T}
\end{Equation}
以上的\xref{eq:内能的微分},\xref{eq:功的微分},\xref{eq:热的微分}具有统一的形式，即可以视为关于状态$(P,V,T)$的一组三元函数$f_P,f_V,f_T$形式写出的$f_P\dd{P}+f_V\dd{V}+f_T\dd{T}$，然而，尽管其具有全微分的形式，但在实质上这未必真的是某个三元函数$F(P,V,T)$的全微分$\dd{F}$，这需要满足以下极为苛刻的条件
\begin{Equation}
    \pdv{f_P}{V}=\pdv{f_V}{P}\qquad
    \pdv{f_V}{T}=\pdv{f_T}{V}\qquad
    \pdv{f_T}{P}=\pdv{f_P}{T}
\end{Equation}
而\xref{eq:功的微分}和\xref{eq:热的微分}中包含了$P\dd{V}$，两者无法满足上式第一项的要求，两者的表达式也就并不是某个三元函数的全微分了，即并不存在所谓的“热$Q(P,V,T)$”和“功$W(P,V,T)$”的状态函数，这与我们理解中“热和功是过程量而不是状态量”的认知是相符的。然而，当我们在表达一些微元关系时，确实存在需要表述“微小的热”和“微小的功”的情景，为此我们就将那些具有全微分形式但实质上不构成全微分的表达式，称为\uwave{不精确微分}（inexact differential）并使用$\fdd{Q},\fdd{W}$的形式表达，对应则以\uwave{精确微分}（exact differential）称呼诸如$\dd{E}$的全微分。
\begin{OldBoxFormula}[热量的第一微分关系]
    热量与状态量$(P,V,T)$具有以下微分关系
    \begin{Equation}&[]
        \fdd{Q}=P\dd{V}+\frac{i}{2}\nu R\dd{T}
    \end{Equation}
\end{OldBoxFormula}
热力学第一定律是能量守恒定律的直接推论，在历史上曾经有不少人企图制造一种，既不消耗系统内能，又无需外界提供能量，而永远对外做功的机器，这类违背热力学第一定律的机器被统称为“\uwave{第一类永动机}”，这样的企图经过无数次实验都已失败告终，因为其试图凭空产生额外的能量。因此，我们其实也可也将热力学第一定律表述为：\empx{第一类永动机是不可能实现的}。

热力学第一定律是关于能量的定律，因此该定律的适用与准静态过程无关，即其可以适用于两个平衡态之间的一切复杂过程，而且对气体、液体、固体等一切物质系统也均是成立的。


